\documentclass[11pt, a4paper]{article}
% ---- Core Geometry & Margins ----
\usepackage[margin=1in, top=1.2in, bottom=1.2in]{geometry}
\setlength{\parindent}{0pt}
\setlength{\parskip}{0.8em}
% ---- Font & Micro-Typography ----
\usepackage[utf8]{inputenc}
\usepackage[T1]{fontenc}
\usepackage{lmodern}
% \usepackage{newtxtext}
% \usepackage{newtxmath}
% \usepackage{libertinus}
\usepackage{microtype}

% ---- Language & Math ----
\usepackage[english]{babel}
\usepackage{amsmath, amssymb, amsthm}
\usepackage{booktabs}  % In your preamble

% ---- Formatting & Aesthetics ----
\usepackage{booktabs}
\usepackage{graphicx}
\usepackage{xcolor}
\usepackage{csquotes}

% ---- Color Palette ----
\definecolor{navyblue}{RGB}{16, 55, 87}
\definecolor{darkslate}{RGB}{53, 64, 83}

% ---- Hyperlinks Setup ----
\usepackage[
  pdfencoding=auto,
  colorlinks=true,
  linkcolor=navyblue,
  citecolor=navyblue,
  urlcolor=navyblue
]{hyperref}

% ---- Bibliography Setup ----
\usepackage[
  backend=biber,
  style=numeric-comp,
  sorting=nyt,
  giveninits=true
]{biblatex}
\addbibresource{references.bib}

% --- Metadata ---
\title{Random Thoughts}
\author{Guetchuin Tadouanla}
\date{today}

% ====================================================================================
\begin{document}
\pagenumbering{roman}
\setcounter{page}{1}

\tableofcontents
\newpage
\pagenumbering{arabic}
\setcounter{page}{1}

\section{Case study "Quality of Product" }
\textbf{Year: Sunday January 1, 1961}
\subsection{Facts}
For many years the ABC Company has manufactured a product which enjoys a high quality rating in
the industry and among the public. Competing manufacturers have now introduced a similar product
of lower quality at lower cost, and this competition has caused a serious decline in the sales
of the product manufactured by the ABC Company. To meet this competition the ABC Company
instructs its engineers to redesign its product in order that it may be made available to the
market at lower cost. Upon receiving these instructions some of the engineers question whether
such an action would be consistent with the Canons of Ethics because a lower quality product
under the same brand name would mislead the public into accepting a product of lesser quality
in the mistaken belief that it meets the high quality standards with which the product has been
associated in the public mind for many years.
\begin{itemize}
  \item  Use to manufacture a High quality product.
  \item Introduction of similar product from competing manufacturers sold at a lower cost.
  \item Decline in sales of the product from ABC company.
  \item ABC company instructs its engineers to redesign the product to lower the cost of the
        product on the market. \item Engineers question the decision in regards to the Canons of Ethics as this seems
        misleading to the public into accepting a product of lesser quality as the name of the
        product remains the same.
\end{itemize}

\subsection{Question}
\textbf{Do the engineers have a proper interest and ethical obligation to protest the company's decision, or to refuse to design a lower quality product?}

\subsection{Personal opinion on the Dilemma}
Keeping the same product name for the lower quality product, could be deciteful to the
public if the company fails at being transparent on the drop in quality to the public. So
servicing the redesign of the product  as a professional Engineer is indirectly
contributing at misleading the public. Though the unfavorable reactions will be directed against the company rather than the Engineer.

\subsection{Objective argumentation}
\begin{itemize}
  \item Canons rules is to protect and serve the public safety, health and welfare by the
        application of engineering knowledge and to protect and defend the high standing of the
        engineering professions. Meaning if the engineer believes that the redesign into a lower
        quality product jeopardizes safety or health, Canon 4 clearly indicates their
        ethical \textbf{responsibility/obligation} to advise management.

  \item Neither the canons nor the Rules, however, deal with business decisions of companies.
        The canons and rules should be restricted to situations which clearly raise questions of
        the ethical conduct of individual engineers.

  \item The lower quality product is exclusively a matter for management determination, based on
        the company's evaluation of its commercial operations. while the Engineers are
        respopndible for giving management honest professional advice.
  \item Neither the Canons nor the Rules, however, deal with business under the Canons or Rules. They may feel that the decision is not a wise one, that it may reflect adversely upon the public acceptance of the Company's product, or even that the public may be misled. These are risks for management evaluation, but there is no reason for the engineers to refrain from offering the company their opinions and comments through normal communication channels within the company. If the public is misled as to the product's quality, the unfavorable reaction will be directed against the company. It would be a disservice to the engineering profession to attempt to extend the Canons and Rules to the point of interference in business decisions of the commercial world.


  \item Neither the Canons nor the Rules, however, deal with business under the Canons or Rules.
        They may feel that the decision is not a wise one, that it may reflect adversely upon the
        public acceptance of the Company's product, or even that the public may be misled. These
        are risks for management evaluation, but there is no reason for the engineers to refrain
        from offering the company their opinions and comments through normal communication
        channels within the company. If the public is misled as to the product's quality, the
        unfavorable reaction will be directed against the company. It would be a disservice to the
        engineering profession to attempt to extend the Canons and Rules to the point of
        interference in business decisions of the commercial world.
\end{itemize}

\section{speaker notes for the presentation}
\subsection{Introduction}
\begin{itemize}
  \item Good morning everyone and welcome to a short presentation on Engineering Ethics on the case
        study of  \textbf{Quality of product}
  \item During this presentation, I will like to share with you an in depth examination of ethics,
        professional responsibiliites
        personal perspectives, and business decision  in engineering.
  \item This is a pretty old case study dating from the 1960s but a really interesting one to me as it
        shows a possible dilemma in which one can potentially find themseleves in as an Engineer.
\end{itemize}

\subsection{Background and Facts}
\begin{itemize}
  \item So let's quickly have a look into the case study itself.

  \item Basically there is this company named  \textbf{ABC}. They have long made high quality products
        and have earned a strong reputation in the industry

  \item Competitors have  introduced a cheaper, low quality version, and ABC's sales have dropped
        drastically.

  \item Management has instructed the Engineering team to produce a new cheaper and low quality design
        of the same product to be able to compete with the current market.

\end{itemize}
\subsection{Concern}
\begin{itemize}
  \item Engineers questioned if retaining the same brand name for a lesser-quality product violates
        the canons of Ethics by misleading the public regarding expected product standards.
\end{itemize}
\subsection{The Core Ethical Question}
\begin{itemize}
  \item Do the engineers have a proper interest and ethical obligation to protest the company's
        decision, or to refuse to design a lower quality product?
\end{itemize}
\subsection{Personal Opinion on the Dilemma}
On the company's point of view:
\begin{itemize}
  \item \textbf{Risk of Deceit}: What I mean by that is, Upon release of the product, consumers or the
        public should me made aware of the drop in quality of the product as a reason for the price
        drop.
  \item
\end{itemize}
On the engineer's point of view:
\begin{itemize}
  \item \textbf{indirect involvement}: Right, like in the case where the Engineer is aware of the lack
        of transparency and still designs the product they are indirectly contributing to misleading
        the public.
\end{itemize}

\textbf{Direction of Fallout}: Since no consumer really knows about the Engineers involved in
designing the product.

\subsection{Ethical Objective Argumentation}
Looking at the Dilemma through the lense of  Etchical Objective Argumentation,....
\end{document}
